\chapter{Discussion}
\label{c:discussion}

\noindent
This chapter interprets the results with reference to the three research questions. Results that
cut across the individual hypotheses are drawn out as stand-alone findings, and the chapter closes
with the limitations of the design and what future research might address.

\section{Interpretation with Respect to the Research Questions}
\label{disc:rqs}

\subsection{RQ1: Does Perceived Authorship Influence Human Emotion Judgments?}
\label{disc:rq1}

\noindent
The extent to which a framing effect is found varies by category.

Human judgments shift significantly under author framing for author-relevant ambiguous texts on
both persona sides, but not for unambiguous control texts on either side. This pattern is in line
with H1. The controls' non-rejection does not prove a zero effect, or a separation of the
persona's content from the other alterations the author cue makes.

The author-independent category is significant under persona~A and not under persona~B. A split
decision of this kind at the significance threshold cannot on its own show that the personas have different effects.
The exploratory direct comparison of their text-level raw TV shifts is 4.44 percentage points,
with a 95\% bootstrap interval from $-2.08$ to $11.06$ points (Section~\ref{res:h1}). Because this
interval includes zero, these data do not establish an asymmetry between A and B. Whether the
category responds to incidental persona content remains a question for further investigation, not a
demonstrated result.

The size of these effects requires a qualification, developed in Section~\ref{res:h3null}. Taking
the vote distribution of each text alone, sampling noise from three votes would account for a
substantial portion of the raw human shift. The effect is detected in the sense that it exceeds
the permutation null, but the raw TV magnitude is a crude figure that reflects not only the effect
of framing but also ordinary sampling noise, so it should not be interpreted as a pure measure of
framing's own contribution to the movement.

The second qualification concerns how the effect was detected in the first place.
Repeating the identical permutation test on the selected label rather than the vote vector rejects
the null in one of six cases under a reverse-alphabetical tie-break convention, and none under the
alphabetical convention used throughout this thesis (only these two conventions were tested;
Section~\ref{res:h1perm}). Thus, nearly all of the support for RQ1 comes from the distributional
representation; under the majority-label representation, one of six tests would have found a
significant effect, and under this thesis's own convention, none. The gap is largely an artefact
of the representation. Its exact size, however, is not intrinsic to the data but is determined by
an otherwise-arbitrary tie-breaking decision.

\subsection{RQ2: Does Perceived Authorship Influence LLM Classifications?}
\label{disc:rq2}

\noindent
Yes, and with fewer qualifications than RQ1 requires.

The ensemble effect is significant on both persona sides. The observed label-change rates are
nonzero for all seven models, although per-model significance tests were not performed. The
exploratory test detects no prompt-configuration effect, but nor does it demonstrate equivalence
across formats. The differences observed across models and formats make the overall
result less dependent on any particular setting, subject to the single-generation limitation
below. This extends the finding of Santurkar et al.\ \citeyear{Santurkar2023WhoseOpinions} that
persona conditioning shifts model outputs, from opinion elicitation to emotion classification, and
matches the direction of Zhang et al.\ \citeyear{Zhang2018Persona} on persona-conditioned dialogue.

The generic-human control makes the claim considerably sharper. The neutral-to-generic-author
contrast is approximately 8\% on both sides (7.98\% and 8.01\%), while the generic-author-to-
specific-persona contrast is 16.2\% and 19.3\%. These distances are sensitive both to the presence
of an attribution and to its specific content, but they are not additive components of a total
effect, so they cannot establish what proportion of the overall shift is due to identity
specifically. Without this control, the two could not have been distinguished at all.

The claim has two limits. First, the significance test identifies \emph{coordinated}
redistribution across models; 61.2\% of the cells contribute no variance to this test's null, and
a scenario in which every cell contained exactly one dissenting model would yield $p = 1$ even
though every cell changed. The complementary raw change rate, 21.3\% of model pairs altering their
label, is reported because the two measures answer different questions. Second, each model
produced one generation per condition, so no distinction is made here between a framing response
and ordinary stochastic variation in generation. That would require multiple generations under the
same prompts, which were not gathered.

\subsection{RQ3: Do LLMs Resemble Humans in Direction and Magnitude?}
\label{disc:rq3}

\noindent
Here the answer splits cleanly in two, and that split is the most substantive finding
of the thesis.

\textbf{In direction, models resemble humans weakly but reliably.} Model delta vectors, when
compared with human ones, have a mean cosine of 0.112, well above a shuffled-pairing null. Models
do not merely shift when framing is applied; they shift in a direction associated with the
direction the same text moves for human readers. A cosine of 0.11 is nonetheless a small
association: the resemblance is genuine but incomplete, as H4 suggested. The criterion is also
more stringent than the usual label-agreement comparison. It does not ask whether the model
reaches the same answer as a human, but whether it \emph{responds to context} in the same
direction, which is closer to the question posed by ambiguity-aware approaches such as Hong et
al.\ \citeyear{Hong2024AERLLM}.

\textbf{In magnitude, models and humans differ substantially.} Humans shift more, by roughly 34
percentage points, and this survives standardising both sides to an equal number of votes: doing
so reduces the estimated gap by only about 1.5 points (roughly 4\%), and it does not remove the
between-person/within-model design asymmetry discussed below.

One further pattern escapes magnitude altogether, although it calls for caution. Under framing,
humans become measurably \emph{less in agreement with one another}: entropy rises significantly in
all three categories. No dispersion effect is detected for the model ensemble in any category.
These are separate tests, and a significant result on one side alongside a non-significant result
on the other does not by itself establish that humans and models differ from each other
statistically; that difference was not tested directly. Taken as two separate results rather
than a demonstrated contrast, the pattern suggests that models redirect their probability mass
while human judgments also become more dispersed. An analysis reporting shift magnitude alone
would never see this.

For the models, then, the results indicate sensitivity to author framing, weak directional
alignment with human judgments, and smaller observed shifts. Human entropy increases were
statistically detected, but not model entropy increases; whether those entropy responses differ
between humans and models remains untested.

\section{Cross-Cutting Findings}
\label{disc:crosscutting}

\subsection{What Aggregation Costs}
\label{disc:aggregation}

\noindent
This finding is a measurement rather than an argument, and it is the clearer of the two
methodological results in this section. A selected-label analysis of these same data would have
recorded 38.7\% of human comparisons and 41.4\% of model comparisons as showing no change, when
the underlying distribution demonstrably moved (under the alphabetical tie-break convention;
crediting the label method for tracking the complete set of tied top emotions instead lowers these
to 25.8\% and 38.7\% respectively). For the model ensemble, the comparisons such an analysis
cannot see outnumber those it detects by more than two to one, and for humans the invisible
comparisons are nearly three-quarters as numerous as the detected changes (38.7\% invisible
against 52.8\% detected).

Label-based analysis is not always inappropriate. For the question asked here, however, whether
framing shifts interpretation, the choice of representation decides whether the phenomenon can be observed at
all. The human side shows this most clearly: a study designed identically but analysed on
majority labels would have detected a statistically significant human framing effect in one of
the six H1 tests (reverse-alphabetical tie-break convention) or none at all (this thesis's
alphabetical convention) rather than three, and would most plausibly have been interpreted as
providing little or no evidence of human sensitivity to author framing. The model-side
selected-label test, by contrast, still detects a significant effect under either representation
(Section~\ref{res:h2magnitude}), so for the model ensemble the choice of representation
changes how much of the shift is visible, not whether an effect is detected at all.

This concern was raised in principle by Aroyo and Welty \citeyear{Aroyo2015CrowdTruth} and Basile
\citeyear{Basile2020GoldStandard}; the figures above put a size on it. They are also consistent with
Pavlick and Kwiatkowski's \citeyear{Pavlick2019InherentDisagreement} argument that disagreement in
subjective annotation is frequently inherent rather than erroneous. That literature argues that
aggregation loses signal; this measurement shows how much is lost and of what kind.

The finding on baseline agreement points the same way. Overall, the human and model dominant
labels agree for only 58.0\% of comparisons under neutral framing (as high as 74.2\% on
unambiguous texts, and as low as about 50\% on the two ambiguous categories), but within both the
matching and non-matching buckets, full-vector similarity spans almost the entire possible range.
Reporting only the binary match rate compresses a continuous and highly variable quantity into a
single number that conceals its own variance, and even that single number shows the two sides
starting from a noticeably different reading of the ambiguous texts before any persona is
introduced.

\subsection{Small Samples Make TV Deceptive}
\label{disc:smallsample}

\noindent
One finding came out of methodological scrutiny rather than the original design, and it applies
well beyond this thesis.

Total variation distance between two small vote samples is not zero when no effect is present, and
it grows with the dispersion of the underlying distribution. Two independent three-vote draws from
the same distribution routinely disagree, and they disagree more when that distribution is spread
out. Since categories defined by ambiguity are, by construction, categories defined by dispersion,
comparing raw TV across such categories confounds the quantity of interest with the measurement's
own noise floor.

Referencing each text against its own conditional sampling null showed that expected null TV is
itself ordered by category, ranging from roughly 43\% to 55\%. The excess above that
conditional-null expectation is on the order of two to six percentage points, not forty-five to
sixty. Any study that compares distributional distance across groups differing in baseline
dispersion runs into this problem, and it stays hidden unless the null is constructed explicitly.

\section{Limitations}
\label{disc:limitations}

\noindent
These results are constrained by the limitations set out below. They have been gathered here
rather than scattered throughout the results chapter, so that they can be weighed together.

\textbf{Vote counts are small and unequal between sides.} Human vectors rest on three raters per
condition and model vectors on up to seven. This gives rise to several structural rather than
behavioural human--model asymmetries reported in Chapter~\ref{c:result} (tie
frequency, new-emotion rates, zero-shift rates, dominant-set stability). Count standardisation
(Section~\ref{res:h4counts}) addresses the magnitude comparison specifically, but not these
descriptive asymmetries.

\textbf{The human comparison is between-person; the model comparison is within-model.} Different
people rated each framing condition, whereas the same model answered both. Part of the measured
human shift is between-rater variability with no model counterpart. The problem is built into the
design, and matching vote counts does not fix it. It is the main caveat on the H4 magnitude
result.

\textbf{H3's categories were partly defined by the property H3 tests.} Texts were assigned to the
author-relevant group using an LLM judge's own sensitivity to whether author framing flips the
reading. H3 is therefore a transfer comparison: whether a grouping derived from that judge's
sensitivity predicts effects in newly collected data, and not an independent discovery that
author-relevance is the operative principle. This is a sharper limitation for the model-side
ordering specifically: the model ensemble's framing effect is reported as ordered by the same
category that the LLM was used to define, so part of that ordering is true by construction rather
than an independent finding.

\textbf{The category ordering is not separable from baseline dispersion.} H3's contrasts are
significant in raw TV, but not significant once referenced against each text's sampling null. This
is not because the null-referenced intervals are wider; they are, if anything, narrower in
absolute width than the raw contrasts' intervals. Rather, the estimated excess over the
conditional null shrinks much more than the interval does, so the same interval now spans zero.
This does not demonstrate equivalence between categories, but it does show that this design cannot
attribute the raw ordering to framing rather than to baseline dispersion at 100 texts per category.

\textbf{Model stochasticity is unmeasured.} One generation per condition means the model figures
are single draws, not stable per-model expectations.

\textbf{No generic-human condition was collected for humans.} The decomposition separating author
presence from persona content is available for models only, so the equivalent human question
remains open.

\textbf{Per-emotion effects are underpowered.} Only 7 of the 66 per-emotion tests survive FDR
correction, and 32 of the 33 human tests do not, including the numerically largest movements in the
heatmap. The broader descriptive pattern may be real, but this study establishes only the one
human effect and the six model effects that survived correction, not the pattern as a whole.

\textbf{Cross-text respondent dependence cannot be modelled, and within-text independence is not
fully guaranteed either.} No collection form repeated a text across both of its conditions, which
is a property of the form design and can be verified from the assignment records. Respondent
identifiers were not retained, however, and there is no way to determine whether the same
individual responded to both conditions of the same text; that is, the same individual could in
principle have rated both conditions without this being evident in the retained data. The same
lack of participant tracking means cross-text correlation, one person's answers on one text being
linked to their answers on another, also cannot be ruled out or quantified: each form carried on
the order of thirty items, and a one-to-one correspondence between forms and participants would
bound the possible correlation structure, but no such correspondence is verifiable from what was
retained. The within-text permutation null also requires comparable respondent groups under the
null; disjoint form assignments narrow this risk but do not, on their own, establish that
exchangeability assumption.

\section{Future Directions}
\label{disc:future}

\noindent
The limitations above suggest five extensions.

\textbf{Collect human annotations at matched vote counts.} If human raters per condition were
raised to seven, the main structural asymmetry would be removed, and the descriptive comparisons
(tie rates, new-emotion rates, dispersion) could be interpreted behaviourally rather than
arithmetically.

\textbf{Add the generic-human condition to the human study.} This would permit the same
neutral-to-generic-author versus generic-to-specific-persona comparison already run on the model
side, so that the two separately measured distances could be compared directly for humans instead
of assumed from the model pattern.

\textbf{Sample repeated generations per model.} Several generations per prompt would separate
framing response from generation variance and convert the model figures from single draws into
estimates with their own uncertainty.

\textbf{Define ambiguity categories independently of the tested property.} If texts were assigned
to categories on criteria not derived from an LLM judge's own framing sensitivity, H3 would become
a direct test rather than a transfer comparison.

\textbf{Retain respondent identifiers.} Anonymised respondent identifiers linked to text and
condition would allow a participant-aware analysis and remove the one dependence structure this
study cannot quantify.

Another direction comes from the entropy finding rather than from a limitation. A significant
human dispersion increase was detected and no model dispersion increase was. The apparent
contrast was not tested directly, and to the author's knowledge the existing literature on persona
prompting has not examined it. It relates to work modelling annotators individually rather than in
aggregate \cite{Davani2022Disagreement, Uma2021Survey}: if human interpretive uncertainty is
itself structured, a model that reproduces answers without reproducing uncertainty is capturing
only part of the behaviour. Can models be made to show human-like interpretive uncertainty, and
not just human-like answers? That question looks both open and tractable.
